\title{Adaptive Reading Budgets for Joint Forecast Audits with Cumulative Data}
\author{Jishan Kapoor\\Independent Researcher\\Toronto M4Y 1C9, Canada\\
Corresponding author: \href{mailto:kapoorjishan19@gmail.com}{kapoorjishan19@gmail.com}}
\date{}
\begin{document}
\maketitle
\begin{abstract}
\textbf{Background.} Model audits share verification readings across
predictors. We ask when choosing later readings after earlier answers
reduces the number needed for precise comparisons.

\textbf{Methods.} We combine sharp comparison intervals, adaptive query
planning and exact randomized-batch certificates, with replay on cumulative meters.

\textbf{Results.} Under compact convex uncertainty, for one comparison,
if a deterministic policy guarantees
precision with two exact scalar readings, some randomized fixed batch succeeds with
worst-case probability at least $1/2$. This lower bound disappears with two
comparisons. For every $m\geq3$, we construct instances where adaptation
always succeeds but the best randomized batch succeeds with worst-case
probability exactly $1/m$. All coordinates are queryable, the budget remains
two, and each comparison is separately certifiable by a fixed batch.
The separation also holds for all pairwise comparisons among three models.
On 225 of 1,098 meter blocks across six systems, two adaptive readings
guarantee joint half-width precision where three fixed readings are necessary.
In 191 cases, every scalar comparison separately needs at most two fixed
readings, isolating a shared-budget obstruction. We certify optimal randomized-batch
precision on all 1,098 blocks. Optional-reading reductions are 6.40\%
and 7.67\% across the panels. Absolute-precision and capacity analyses
identify when simpler fixed plans suffice. Unrestricted bulk retrieval
is faster on the tested public interface.

\textbf{Conclusions.} Sharing a budget can make feedback necessary even when
separate comparisons need none. The meter study identifies when feedback
saves readings and when simpler protocols suffice.
\end{abstract}
\noindent\textbf{Keywords:} artificial intelligence; machine learning;
adaptive information; minimax risk; optimal recovery; robust optimization.

\section{Introduction}
A model audit may need more verification data before it can report a precise
comparison. The auditor can retrieve a complete record, choose a batch of
readings, or choose later readings after seeing earlier answers. We study
this last option, adaptive querying. Our output is a sharp interval: the
smallest interval containing the loss difference for every record compatible
with the observations. Its width measures how much the reported comparison
can still change. We ask when feedback is needed to meet a stated width
requirement with the same reading budget.

This is an acquisition problem for fixed predictors. The auditor specifies
which loss contrasts must be precise, what information is already available,
and which measurements may be requested. We first compute the smallest
sufficient fixed reading budget, then test whether two adaptive readings
can reduce it. The algorithm returns a reading plan and checks its intervals.
Identifying a better model is a separate
condition: a comparison interval must also exclude zero.

The central result is a boundary between one comparison and a shared
collection of comparisons. For one affine contrast, a deterministic policy
that guarantees precision with two exact scalar readings implies a
randomized fixed batch with worst-case success at least one half
(Theorem~\ref{thm:helly-precision}). With two contrasts, that floor disappears.
For every $m\geq3$, Theorem~\ref{thm:two-target-collapse} constructs instances
with optimal adaptive success $1$ and optimal batch success $1/m$, still
using two readings. Each contrast separately has a successful fixed batch;
all coordinates are queryable at equal cost. The same separation holds for
all pairs among three nonnegative forecast vectors
(Corollary~\ref{cor:all-pairs-collapse}). The dictionary grows with $m$,
and three fixed queries suffice. Moreover, the optimal batch mixture is
unique: it assigns mass $1/m$ to each of $m$ necessary batches. Even positive
worst-record success requires all $m$ batches. This exact support requirement
isolates what fails when separately sufficient measurements share a budget.

We make that distinction computable for cumulative-meter audits. An exact
two-query planner, minimum fixed-budget calculation and matching randomized-batch
certificates determine when feedback saves a reading. The study uses frozen
forecasts on 1,098 blocks from six systems. Separate scalar-budget calculations
isolate the shared-budget mechanism; absolute precision, nearby thresholds
and capacity controls test its sensitivity on the same records. Replay then
measures what the policies achieve on stored answers. A public-API control
distinguishes the mathematical reading budget from the available retrieval
interface.

The same calculation can rule out an insufficient budget. Explicit adverse
records show that randomizing an adaptive policy cannot restore certainty
on any of the 583 blocks without a deterministic guarantee
(Proposition~\ref{prop:assurance-feasibility}). Thus the method supports a
three-way decision: choose a fixed batch, use feedback, or obtain more
information.

The paper develops the scalar bound, the shared-budget construction and the
meter planner, then reports its certificates, recorded outcomes and controls.
The retrospective study reuses frozen source panels. Online Resource~1
retains the supporting scalar penalty classification (S11), a separate
mean-joint-width construction (S14), historical controls and executable
source-to-table verification.

\subsection{Related work and the contribution boundary}
Partial identification studies conclusions that survive incomplete
observations \citep{manski,polo}. Cancellation in paired squared losses
already yields an affine function of the missing truth \citep{patton}.
Active evaluation and model selection acquire labels to estimate performance
or choose predictors \citep{sawade,hara,active,okanovic2025}. We instead
require the entire compatible interval of each fixed-record contrast to
meet a declared precision target. \citet[Theorem~1 and Proposition~1]{kuboyama2026}
studies completion-based model-choice certificates with adaptive binary
labels and shared paired bounds. Its distinction between hindsight certificate
size and acquisition requirements (Theorem~2 and Proposition~3) is also
relevant here; our guarantees concern joint widths for continuous targets.

The width of an information fiber is the established local diameter of
information \citep[Section~4.1]{novak2008}. Our shared objective is its
weighted $\ell_\infty$ version for a vector target. Classical deterministic
adaptivity bounds cover linear information on convex symmetric classes
\citep[Theorems~4.4--4.5]{novak2008}; nonsymmetric recovery is studied by
\citet{novak1995}, and noisy scalar criteria by
\citet[Theorems~2--3 and~5]{osipenko}. Modern randomized recovery minimizes
reconstruction error \citep[Section~2 and Proposition~4.1]{krieg2025}
\citep[Section~9.1, Eq.~(9.15), author manuscript v2]{krieg2026}.
Adaptive vector reconstruction also admits high-probability guarantees and
growing gains over nonadaptive reconstruction
\citep[Theorems~3.1 and~4.1, author manuscript v1]{kunsch2026}.
Our precision objective requires the full terminal fiber to be narrow.
The shared-budget construction keeps the output dimension and hard query
budget at two as the predeclared coordinate dictionary grows.
For supporting mean-width results, the order of the
worst-record and randomization operations also differs from expected global
radius \citep[Eq.~(2.4), Remark~6.1, author manuscript v2]{krieg2026}.
Online Resource~1, S11 gives that comparison and the scalar penalty
classification, whose sufficient direction applies Sion's minimax theorem
\citep[Theorem~3.4]{sion1958} and Jensen's inequality.

Information-dependent decisions are established in adjustable robust
optimization \citep{wei2024}. \citet[Section~3.2, Eq.~(3) and
Theorem~3.2, author manuscript v3]{ddid} give equivalent
observation-consistent recourse and information-fiber formulations.
Our scalar precision bound adapts the
sufficient-policy result of \citet[Theorem~2 and Example~5]{kurtz2026}.
Kurtz's Theorem~11 also permits policy counts to grow with recourse-region
complexity. His Example~20 gives a particularly close construction: a
scalar uncertainty coordinate selects an interval on which a recourse
policy is feasible, and the chosen policies must cover those intervals.
Our construction shares this interval-cover mechanism. The difference is
how it arises: we derive the intervals from the sharp widths of actual
coordinate-observation fibers, with every coordinate available at the same
cost. Each scalar conditional width is affine, but their joint maximum
loses the concavity used by the scalar bound. Two
affine targets and two unit-cost readings have exact success values $1$
and $1/m$, although each target alone has a sufficient fixed batch.
The adverse records further force support $m$ for any positive batch
success and uniquely determine the optimal mixture. The output dimension
and budget stay fixed while the measurement dictionary grows. This is the
specific contribution beyond general policy-count bounds and the established
benefits of adaptive vector recovery.

Adaptive cumulative sampling has close precedents.
\citet[Section~4.1.1, Corollary~1 and Theorem~3]{gaillard2023} studies
certified scalar integral estimation from a monotone function. For monotone
contrast coefficients, summation by parts relates our meter problem to
weighted monotone integration. We pair an exact shared-query planner with
continuum covering and adverse-record certificates, so each reported
precision separation has matching bounds under one measurement interface.
The separate aggregate solver in Online Resource~1, S5/S7 builds on
weighted-oracle randomization \citep{kawase2024}, dynamic-program LP
representations \citep{martin1990}, and refinement of tied robust solutions
\citep{iancu2014}; it is supporting computation rather than the central
shared-query contribution.

\section{Model Comparisons and Sharp Information}
Let $a,b\in\mathbb R^N$ be two fixed prediction vectors for the same
scheduled targets. Their paired mean-squared-error difference is
\begin{equation}
 D(y)=\frac1N\sum_{i=1}^N\{(a_i-y_i)^2-(b_i-y_i)^2\}
 =d+c^\top y,\qquad c_i=\frac{2(b_i-a_i)}N.
 \label{eq:linear}
\end{equation}
All scheduled targets remain in the denominator. Initially observed values,
documented aggregate measurements and other defensible constraints define
the feasible set $F$. Any permitted archive record that directly reconstructs
a target is incorporated before this set is declared. Sharpness is always
relative to that declared information set.

The attainable interval of $D$ on $F$ is an identification interval for
this record, not a sampling confidence interval. A negative upper endpoint
certifies that $a$ has smaller squared loss for every completion; a positive
lower endpoint certifies $b$. An interval containing zero cannot certify
a strict winner. Subsequent exact measurements restrict $F$ to a fiber and
therefore cannot widen its sharp interval. A prescribed width threshold
instead asks for a precision guarantee, regardless of the interval's location.

The mathematical examples below also describe paired squared-loss widths.
Given any contrast coefficient $c$, nonnegative predictions
$a_i=N\max(-c_i,0)/2$ and $b_i=N\max(c_i,0)/2$ realize that coefficient
in \eqref{eq:linear}. The induced constant changes the interval's position
while preserving its width. This gives the constructions a paired
model-comparison interpretation.

